\clearpage
\chapter*{Abstract}

With the deepening of smart manufacturing and Industry 4.0, constructing high-fidelity and real-time virtual replicas of physical entities has become a core requirement for the implementation of industrial digital twin technologies. In this context, virtual-real mapping serves as a key enabling technology connecting the physical and information worlds, and the quality of its underlying 3D visual reconstruction base directly determines the fidelity and usability of the twin system. Recently, explicit neural rendering technologies represented by 3D Gaussian Splatting (3DGS) have brought a paradigm-shifting breakthrough to novel view synthesis with extremely photorealistic rendering and remarkably high forward rasterization frame rates. However, when directly applied to industrial-like digital twin scenes—which typically contain numerous highly reflective metal components, large areas of textureless solid-color casings, and complex equipment occlusions—traditional 3DGS exposes severe adaptability issues. On the one hand, pure-vision photometric constraints easily fail on complex materials, triggering geometric topological collapse and pervasive "floater" artifacts. On the other hand, the disordered growth of massive discrete Gaussian primitives introduces enormous high-dimensional parameter redundancy and extremely high VRAM consumption, making it difficult to meet the large-scale deployment requirements of edge devices and Web browsers. To address these challenges, this thesis aims to achieve efficient and high-fidelity virtual-real mapping for industrial-like twin scenes by systematically investigating the spatial geometric constraint enhancement and lightweight model compression of 3DGS. The main research work and innovative contributions are as follows:

\textbf{First, a multi-modal adaptive geometric constraint reconstruction algorithm for industrial-like scenes is proposed.} To tackle the ill-posed reconstruction problems caused by strong anisotropic reflections and textureless regions on industrial equipment surfaces, this thesis breaks the limitation of pure-vision optimization. Initially, through hardware-level spatiotemporal alignment, the absolute physical scale dense point clouds acquired by LiDAR and the relative depth fields inferred by a monocular visual foundation model are cross-modally registered, injecting a robust 3D physical initial skeleton into 3DGS. Subsequently, an adaptive confidence mask mechanism based on depth local high-frequency feature perception is designed to dynamically partition flat trusted regions and detail-exempted regions. Finally, a loss function combining a Huber depth penalty and first-order normal smoothness regularization is constructed to force the Gaussian primitives to adhere to real physical surfaces under the dynamic scheduling of the mask. Experiments show that this algorithm effectively eliminates floater noise on metal surfaces, perfectly repairs geometric tearing on solid-color cabinets, and intactly preserves the high-frequency edges of machining textures while significantly reducing depth errors and Chamfer distances.

\textbf{Second, a lightweight 3DGS representation and edge rendering optimization algorithm for typical twin scenes is proposed.} Targeting the computing power and VRAM bottlenecks caused by the massive parameter volume of 3DGS, this thesis tackles the problem from both data reconstruction and rendering pipelines. At the geometric purification level, a multi-dimensional importance assessment model integrating transmittance, global view visibility, and spatial scale is constructed. Combined with a long-tail distribution adaptive threshold, it accurately strips away massive "ghost Gaussians" hidden in invisible regions. At the attribute dimension reduction level, Vector Quantization (VQ) technology is introduced to map high-dimensional Spherical Harmonics (SH) coefficients, which dominate storage overhead, into discrete codebooks. This is combined with a Straight-Through Estimator (STE) to accomplish Quantization-Aware Training (QAT), avoiding color banding caused by hard quantization. Furthermore, the rendering pipeline adapted for the Web frontend is reconstructed, utilizing WebGPU compute shaders for in-VRAM hardware-level decoding, and combined with the local sorting optimization of mobile TBDR architectures, thoroughly smoothing the real-time fluent roaming path for extremely high-compression-rate models on ultrabooks and mobile devices.

\textbf{Third, an algorithmic pipeline fusing the improved 3DGS is constructed and system-level comprehensive verification is completed.} This thesis deeply integrates "multi-modal geometric constraints" and "lightweight pruning and quantization" into the training pipeline and designs a multi-stage progressive training strategy. The research reveals a strong synergistic effect between the two: the stable physical geometric base provides precise safety guidance for drastic pruning operations. Comprehensive tests on open-source benchmark datasets (Tanks and Temples) and a self-collected complex server room twin dataset demonstrate that, while maintaining an extremely high rendering fidelity of 31.20 dB, the fused algorithm compresses the physical storage volume of the scene model by approximately 95.8\% (down to 48.2 MB), drastically reduces peak rendering VRAM to 0.45 GB, and achieves an ultimate fluent rendering of 120 FPS under standard hardware environments.

In summary, the research results of this thesis effectively break through the material generalization bottlenecks and edge computing power barriers of pure-vision 3DGS in industrial-like scenes. It provides a practical, lightweight, and comprehensive solution for manufacturing enterprises to construct digital twin 3D visual bases with low computing power thresholds, high frame rates, and high fidelity, possessing significant theoretical value and broad engineering application prospects.

\vspace{1em}
\noindent\textbf{Keywords:} Industrial Digital Twin; 3D Gaussian Splatting; Multi-modal Geometric Constraints; Lightweight Compression; Virtual-Real Mapping; Adaptive Rendering